%! Operations with Polynomials — Unit 3, Lesson 3.1 — Slides (Beamer)
\documentclass[aspectratio=169, 11pt]{beamer}
\usepackage{algebra2-beamer}   % provides \forestheader{} and \sectionlabel[color]{}

\begin{document}

% TITLE SLIDE (CourseName is not defined in beamer, so the name is literal)
{
\setbeamercolor{background canvas}{bg=forest}
\begin{frame}[plain]
  \vspace{2.8cm}\hspace{0.55cm}%
  \begin{minipage}{0.9\textwidth}
    {\color{goldacc}\bfseries\small LESSON 3.1}\\[0.5cm]
    {\color{white}\bfseries\LARGE Operations with Polynomials}\\[1.2cm]
    {\color{forestmid}\small Algebra 2~$\cdot$~Unit 3: Polynomial Functions}
  \end{minipage}
\end{frame}
}

% HOOK
\begin{frame}[t, plain]
  \forestheader{Two names for yesterday's curve --- same function?}
  \sectionlabel{Hook}
  \begin{columns}[T]
    \begin{column}{0.52\textwidth}
      \[ g(x) = (x-1)^2(x+2) \]
      \[ g(x) = x^3 - 3x + 2 \]
      \vspace{2pt}
      \begin{itemize}\setlength{\itemsep}{4pt}
        \item Both give $2$ at $x=0$, $4$ at $x=2$, $0$ at $x=-2$.
        \item Three matching points is \emph{evidence}, not \emph{proof}.
      \end{itemize}
    \end{column}
    \begin{column}{0.46\textwidth}
      \begin{block}{What each costume is good for}
        \textbf{Factored} $\Rightarrow$ the \emph{zeros}. \\[3pt]
        \textbf{Standard} $\Rightarrow$ the \emph{degree} and \emph{leading coefficient}, hence the
        \emph{end behavior}.
      \end{block}
      \vspace{3pt}
      {\footnotesize \textbf{Today:} travel from factored to standard by multiplying.}
    \end{column}
  \end{columns}
\end{frame}

% NAMING
\begin{frame}[t, plain]
  \forestheader{Standard form comes first}
  \sectionlabel{Definition}
  \begin{columns}[T]
    \begin{column}{0.55\textwidth}
      Powers from highest to lowest:
      \[ a_nx^n + a_{n-1}x^{n-1} + \dots + a_1x + a_0 \]
      \begin{itemize}\setlength{\itemsep}{4pt}
        \item \textbf{Degree:} the highest power.
        \item \textbf{Leading coefficient:} that term's coefficient.
        \item \textbf{Terms:} 1 monomial, 2 binomial, 3 trinomial.
      \end{itemize}
    \end{column}
    \begin{column}{0.42\textwidth}
      \begin{block}{Two variables}
        A \emph{term's} degree is the \textbf{sum} of its exponents: \\[3pt]
        $-7x^2y^3$ has degree $5$.
      \end{block}
      \vspace{2pt}
      {\footnotesize $-5x^2+7x^4-3+x \;\Rightarrow\; 7x^4-5x^2+x-3$: degree $4$, lead $7$.}
    \end{column}
  \end{columns}
\end{frame}

% ADD & SUBTRACT
\begin{frame}[t, plain]
  \forestheader{Adding \& subtracting: collecting, not creating}
  \sectionlabel[goldacc]{The sign trap}
  \begin{columns}[T]
    \begin{column}{0.5\textwidth}
      \footnotesize
      \textbf{Like terms} = same variables, same exponents. \\
      $3x^2$ and $3x^3$ never combine.
      \vspace{6pt}
      \[ (3x^3-5x+2)+(x^3+4x^2-7) \]
      \[ = 4x^3+4x^2-5x-5 \]
    \end{column}
    \begin{column}{0.5\textwidth}
      \footnotesize
      A minus sign in front is a factor of $-1$ --- it flips \emph{every} sign inside.
      \[ (3x^3-5x+2)-(x^3+4x^2-7) \]
      \[ = 3x^3-5x+2-x^3-4x^2+7 \]
      \[ = 2x^3-4x^2-5x+9 \]
    \end{column}
  \end{columns}
  \vspace{2pt}
  \begin{block}{A degree can drop}
    $(x^4+2x)-(x^4-5)=2x+5$ --- the leading terms cancel, so a sum or difference has degree
    \emph{at most} the larger degree.
  \end{block}
\end{frame}

% MULTIPLY
\begin{frame}[t, plain]
  \forestheader{Multiplying: every term times every term}
  \sectionlabel{One rule, three sizes}
  \begin{columns}[T]
    \begin{column}{0.46\textwidth}
      \footnotesize
      \begin{itemize}\setlength{\itemsep}{5pt}
        \item \textbf{Monomial $\times$ polynomial:} \\ $-2x^2(3x^3-4x+6)=-6x^5+8x^3-12x^2$
        \item \textbf{Binomial $\times$ binomial:} \\ $(2x-5)(3x+4)=6x^2-7x-20$
        \item \textbf{Two variables:} \\ $(3x+2y)(x-4y)=3x^2-10xy-8y^2$
      \end{itemize}
      \vspace{3pt}
      Coefficients \emph{multiply}; exponents \emph{add}.
    \end{column}
    \begin{column}{0.52\textwidth}
      \centering
      {\footnotesize $(x+3)(x^2-2x+4)$ --- six products, no drops:}\\[4pt]
      \renewcommand{\arraystretch}{1.5}
      {\footnotesize
      \begin{tabular}{|c|c|c|c|}
      \hline
      $\times$ & $x^2$ & $-2x$ & $+4$ \\ \hline
      $x$  & $x^3$  & $-2x^2$ & $4x$ \\ \hline
      $+3$ & $3x^2$ & $-6x$   & $12$ \\ \hline
      \end{tabular}}\\[5pt]
      {\footnotesize $= x^3 + x^2 - 2x + 12$}
    \end{column}
  \end{columns}
\end{frame}

% DEGREE OF A PRODUCT
\begin{frame}[t, plain]
  \forestheader{Know the shape before you expand}
  \sectionlabel[goldacc]{Degrees add, leads multiply}
  \[ (2x^3 - 1)(-3x^4 + x) \]
  \begin{itemize}\setlength{\itemsep}{5pt}
    \item Only the two \emph{leading} terms can make the highest power: $2x^3\cdot(-3x^4) = -6x^7$.
    \item \textbf{Degrees add:} $3+4=7$. \quad \textbf{Leading coefficients multiply:} $2\cdot(-3)=-6$.
    \item Degree \emph{odd}, lead \emph{negative} $\Rightarrow$ (Lesson 3.0) left arm \textbf{up},
          right arm \textbf{down}.
  \end{itemize}
  \begin{block}{Why it never fails}
    Nothing can cancel the top term --- there is only one way to build it. So a \emph{product} never
    loses degree, even though a \emph{sum} can.
  \end{block}
\end{frame}

% SPECIAL PRODUCTS
\begin{frame}[t, plain]
  \forestheader{Special products --- and the year's worst error}
  \sectionlabel{Identities}
  \begin{columns}[T]
    \begin{column}{0.52\textwidth}
      \[ (a+b)^2 = a^2 + 2ab + b^2 \]
      \[ (a-b)^2 = a^2 - 2ab + b^2 \]
      \[ (a+b)(a-b) = a^2 - b^2 \]
      \vspace{2pt}
      {\footnotesize $(3x-4)^2 = 9x^2-24x+16$ \\[3pt]
       $(2x+7)(2x-7) = 4x^2-49$ \\[3pt]
       $(5x-3y)^2 = 25x^2-30xy+9y^2$}
    \end{column}
    \begin{column}{0.46\textwidth}
      \begin{block}{$(x+4)^2 \ne x^2+16$}
        Test at $x=1$: \\[3pt]
        $(1+4)^2 = 25$ \quad but \quad $1+16=17$. \\[4pt]
        The missing piece is the middle term $8x$: \\[3pt]
        $(x+4)^2 = x^2+8x+16$.
      \end{block}
    \end{column}
  \end{columns}
\end{frame}

% CLOSING THE LOOP
\begin{frame}[t, plain]
  \forestheader{Closing the loop on the Hook}
  \sectionlabel[goldacc]{$(x-1)^2(x+2)$}
  \begin{columns}[T]
    \begin{column}{0.52\textwidth}
      \centering
      {\footnotesize Step 1: $(x-1)^2 = x^2-2x+1$}\\[5pt]
      \renewcommand{\arraystretch}{1.5}
      {\footnotesize
      \begin{tabular}{|c|c|c|c|}
      \hline
      $\times$ & $x^2$ & $-2x$ & $+1$ \\ \hline
      $x$  & $x^3$  & $-2x^2$ & $x$ \\ \hline
      $+2$ & $2x^2$ & $-4x$   & $2$ \\ \hline
      \end{tabular}}
    \end{column}
    \begin{column}{0.46\textwidth}
      \footnotesize
      \begin{itemize}\setlength{\itemsep}{5pt}
        \item $-2x^2$ and $+2x^2$ \textbf{cancel} --- that is why there is no $x^2$ term.
        \item $x - 4x = -3x$.
        \item Result: $\;x^3 - 3x + 2$. \; \textbf{Proved.}
        \item Degree $3$, lead $1$ $\Rightarrow$ left down, right up --- exactly yesterday's graph.
      \end{itemize}
    \end{column}
  \end{columns}
\end{frame}

% ROADMAP
\begin{frame}[t, plain]
  \forestheader{Where Unit 3 goes next}
  \sectionlabel{Connections}
  \textbf{Today:} add, subtract, and multiply polynomials in one and two variables; special products;
  degree and leading coefficient of a result.\\[8pt]
  \begin{itemize}\setlength{\itemsep}{3pt}
    \item \textbf{3.2} Advanced factoring --- today's products run \emph{backward}, plus the sum and difference of cubes.
    \item \textbf{3.3} Dividing polynomials; Remainder \& Factor Theorems.
    \item \textbf{3.4--3.5} Zeros: Rational Root Theorem; complex zeros \& the FTA.
    \item \textbf{3.6} Graphing polynomial functions \& modeling.
  \end{itemize}
\end{frame}

\end{document}
